\chapter{物理定律的对称性}

\section{对称操作}

我们可以把本章的主题称作为物理定律的对称性 。在矢量分析（第11章），相对论（第16章）以及转动（第20章）等章节中，我们已经讨论过有关的物理定律中对称性的一些特点。

为什么我们要关心对称性呢？首先，在人们的心目中，对称性是非常吸引人的，我们都喜欢具有某种对称性的物体或图案。有趣的是，大自然常常在我们周围所遇到的物体中显示出某种对称性。或许可以想象的最对称的物体是球体，而在自然界中就充满了球体——恒星、行星、云层中的水滴等等。在岩石中找到的晶体也呈现出各种各样的对称性，对它们的研究使我们知道了有关固体结构的某些重要情况。即使动植物世界也显示出某种程度的对称，虽然，一朵花或一只蜜蜂的对称性不像晶体中的对称性那样完美或重要。

但是，我们在这里主要关心的不是自然界的物体 往往是对称的这个事实。我们倒更希望考察宇宙中的一些更引人注目的对称性——即存在于支配物理世界运转的基本定律自身 中的对称性。

首先我们要问，对称性是什么？一条物理定律 怎么会是“对称的”？定义对称性的问题是一件有趣的事情，我们曾提到过外尔给出了一个很好的定义，其要点为，如果有一样东西，我们可以对它做某种事情，在做完之后，这个东西看起来仍旧和先前一样，那它就是对称的。例如，一个对称的花瓶就是这类东西，如果我们使它反射或转动，结果看上去仍旧和先前一样。目前我们要考虑的是，可以对物理现象或实验中的物理状况做些什么事，而其结果却和未做前一样。表\ref{Table:52:1}列举了使种种物理现象得以保持不变的已知操作。

\begin{table}[H]
    \caption{对称操作}
    \label{Table:52:1}
    \centering
    \small
    \begin{tabularx}{0.8\textwidth}{|>{\centering\arraybackslash}X||>{\centering\arraybackslash}X|}
        \hline 
        空间平移 & 空间反射 \\ 
        时间平移 & 全同原子或全同粒子的交换 \\ 
        转过一定的角度 & 量子力学的相位 \\ 
        匀速直线运动变换（洛伦兹变换） & 物质-反物质(电荷共轭) \\ 
        时间反演 & ~ \\ \hline 
    \end{tabularx}
\end{table}

\section{空间与时间的对称性}

我们尝试做的第一件事情，比方说，就是使现象在空间中平移 。如果我们在一定的位置上做一个实验，然后在空间的另一个位置上建立另一套仪器（或者把原来的仪器搬过去），那么凡是在前一套仪器中按一定的时间顺序发生的一切，在后一套仪器中也将以同样方式出现；只要我们安排好同样的条件，并且对前面讲过的一些约束予以应有的注意，即周围环境中所有使仪器不能同样工作的特征都要排除掉。关于在这些情形中怎样确定应包括的因素，我们已经谈过，这里就不再去考虑这些具体细节了。

在今天，我们也同样相信，时间的移动 对物理定律也不会有影响（这是就我们迄今所知而言——所有这些事情都是就我们迄今所知而言）。这意味着，如果我们制造一套仪器，并且在某个时刻，比如星期四上午10:00使它开始工作，然后又制造同样一套仪器，并在（比方说）三天之后在同样的条件下使它开始工作，那么不管何时使这两套仪器工作的情形作为时间的函数是完全相同的。然而我们还是要假设，环境中的有关特征也要及时 作相应的变动。当然这个对称性意味着，如果某人三个月前曾买进通用汽车公司的股票，假如他现在买进这些股票，所遇到的情况将完全一样！

我们也必须注意到地理情况的差别，因为地球表面各处的特征显然是不同的。例如，假如我们在测量了某处的地磁场后将仪器移到另一处去，那么由于地磁场不同，仪器可能不再以完全相同的方式工作，但我们说这是由于磁场与地球有关。我们可以设想，如果使地球和仪器一起移动，那么仪器的工作情况就不会受到影响了。

我们曾相当详细地讨论过的另一件事是在空间的转动：如果把仪器转动一个角度，并且假定其他每件与它有关的物体也随之一齐转动，那么仪器将会同样地工作。事实上，在第11章中，我们曾比较具体地讨论过空间转动中的对称性，并且为了尽可能简洁地处理它，我们还创造了一种称为矢量分析 的数学系统。

在较高级的水平上，有另一种对称性——匀速直线运动的对称性。这个相当不寻常的效应所说的是，如果我们有一件仪器按一定的方式工作，现将该仪器放到一辆汽车里，并使汽车以及与之有关的周围物体都沿直线匀速前进，那么汽车中所出现的物理现象并不会有什么不同：所有的物理定律都显得相同。我们甚至还知道怎样用比较专门性的方式来表示这一点，即在洛伦兹变换 下，物理定律的数学方程式必须不变。事实上，正是在有关相对论问题的研究中，使得物理学家将注意力集中于物理定律的对称性方面。

上面所提的对称性都具有几何的性质，时间与空间多少是类似的，但是还有别的一类对称性。比如，有种对称性就描述了这样的事实：一个原子可以用同一类的另一个原子来替换；换句话说，存在着 同一类的原子。我们可以找到一群原子，如果把其中的一对交换一下，并不会造成什么差别——这些原子是全同的。无论某种类型的一个氧原子会做什么，这类氧的另一个原子也会这样做。有人会说：“真可笑！这正是同一类型的定义 嘛！”这或许只是一个定义，但我们并不知道究竟存在不存在任何“同种类型的原子”，而事实则是，的确存在着许许多多同一类型的原子。因此，当我们说如果用同一类型的一个原子替换另一个而不会出现什么差别时，确实是有意义的。从上述意义来说，构成原子的那些所谓基本粒子也是些全同粒子——所有电子都相同，所有质子都相同，所有正$\pi$介子都相同，等等。

在列举了这么多使现象不改变的操作后，人们或许认为我们实际上能做任何事了；那么就让我们举一些反例，以便看出情况的差别。假设提出这样一个问题：“尺度改变了，物理定律是否对称？”设想我们先制造一台仪器，再制造一台每个部件都放大五倍的仪器，那么它们是否会同样精确地工作？在这种情况下，答案是，不会！例如，从一个装有钠原子的装置中发射出的钠原子光波的波长，与另一个体积为其五倍的装置所发射的钠原子光波波长相比，后者并非前者的五倍，而实际上完全相同。可见，波长与发射装置的大小之比值将会改变。

另一个例子是：在报纸上，有时我们看到由小火柴棒搭成的大教堂的照片，这些惊人的艺术作品是一些退休者用火柴棒粘成的，它比任何真实教堂都精致和奇特。如果我们设想这个用火柴棒制成的教堂果真按真正的尺寸建造起来，就会看到麻烦何在了，它不可能存在下去，整座教堂都会倒塌。因为按比例放大的火柴棒根本不够牢固。或许有人会说：“不错。但是，我们也知道，当存在一种外界的影响时，它也必须按比例改变！”这里谈的是物体承受万有引力的能力。因此我们应先得取一个真正的用火柴棒制成的教堂模型以及真实的地球，我们知道在这种情况下，“教堂”是牢固的。然后，我们该有一个较大的教堂和一个较大的地球，但是这样一来情况变得更坏，因为万有引力增加得更快！

当然，今天我们是根据自然界中的物质由原子构成这一点出发来理解现象与尺度有关这个事实的，很明显，如果我们制造一个小到比方说其中只有五个原子的仪器，那么这种东西肯定不可以任意地放大或缩小。单个原子的尺寸根本不是任意的，而是完全确定的。

物理定律在尺度变化下并不保持不变这一事实是伽利略发现的。他认识到材料的强度并不恰好与其尺寸成比例，同时用表示两根骨骼的图画来说明那个我们刚刚提到的用火柴棒搭成的教堂的问题上讨论过的性质，图中有一根是通常的狗骨骼，它与支撑的狗的重量成适当的比例，而另一根是假想的“超级狗”——比方说大十倍或一百倍的狗的骨骼，这根骨骼按完全不同的比例画得又大又结实。伽利略是否曾把论证真正引申到这样一个结论：大自然的定律必须具有一定的尺度，我们不得而知。但是上述发现给他留下了深刻的印象，以致他认为这件事与发现运动定律同样重要，因为他把这两件事一起发表在同一本名为《论两种新科学》的书中。

另外一个物理定律不对称的例子是我们熟知的：当一个系统作均匀的角速度转动时，其中的表观物理定律与一个不转动的系统的物理定律显得不相同。如果我们安排好一个实验，然后把所有的东西放到一个宇宙飞船里，再在宇宙空间中让飞船本身以恒定的角速度自转，实验仪器将不会照原样那样工作，因为我们知道，在飞船里的东西会被甩出去，并且会发生其他一些情况。这是由于离心力或科里奥利力等而造成的。事实上，我们不必向外看，只要利用傅科摆就可觉察到地球在旋转。

下面我们来叙述一个很有趣的对称性，即时间的可逆性 。初看起来，这显然不成立。很明显，物理定律在时间上是不可逆的，因为我们知道，所有明显的现象在大尺度上都是不可逆的：“挥笔写字，写完再写，……”到现在为止我们所能说的是，这种不可逆性是由于所牵涉的粒子的数量极其巨大而产生的，倘若我们能够看到单个分子，就将无法辨别变化是往正方向发展还是往逆方向发展。更确切地说：我们先制造一台小小的仪器，就能知道其中所有原子的行为，也能观察到它们的运动。然后再制造一台类似的仪器，这台仪器开始工作时的状况与前一台仪器的最终状况相同，但所有的原子的速度正好相反。那么，这台仪器将经历完全相反的变化过程 。换句话说，如果我们拍一部影片，详细地记录了一块材料的所有内部情况，然后再倒过来放，没有一个物理学家会说：“这是违反物理定律的，有些地方搞错了！”当然，如果我们没有去观察所有的细节，事情将是完全明确的，比如说，当我们看见一个鸡蛋落在人行道上，使蛋壳破碎时，肯定会说：“这是不可逆的，因为如果把这件事拍成影片，然后倒过来放时，破碎的蛋壳将会重新拼合，成为完整的鸡蛋，这显然是荒谬的！”但是，如果我们观察单个原子本身，定律看来完全是可逆的。当然，发现这一点要难得多，但很清楚，在微观的基本的水平上，物理学的基本定律在时间上确实是完全可逆的。

\section{对称性与守恒定律}

至此，有关物理定律的对称性已显得十分有趣，但是结果发现，在量子力学中，这种对称性变得更为有趣和更令人兴奋。有这么一件事实：在量子力学中，对于每一个对称的规律都有一条守恒定律与之相对应 。这个最深奥和最美妙的事实对许多物理学家来说简直令他们感到震惊。鉴于我们现在的讨论水平，我们无法对之作更多的说明。物理定律的对称性与守恒定律之间存在着一定的联系。在这里我们只加叙述而不打算作任何解释。

举例说，物理定律对空间平移是对称的。如果与量子力学的原则相结合，结果就意味着动量是守恒的。

物理定律对时间平移是对称的。在量子力学中就意味着能量是守恒的。

关于空间转动一定角度后的不变性与角动量守恒定律 相对应。这些关系是非常有趣和非常美妙的！它们堪称物理学中无比优美和意义深远的东西。

顺便提一下，量子力学中出现的有一些对称性并没有经典的类比，无法以经典物理的方式描述。其中有一个就是：如果$\psi$ 是某个过程的概率波幅，我们知道$\psi$ 的绝对值的平方就是这个过程出现的概率。现在，如果有人进行计算时不用这个$\psi$ ，而是用另一个$\psi'$，它与$\psi$ 只是相差一个相位因子（令$\Delta$为某个常数，把$e^{i\Delta}$ 乘以原来的$\psi$ 即得$\psi'$），那么，作为该事件概率的$\psi'$的绝对值平方就等于$\psi$的绝对值平方
\begin{equation*}
    \psi' = \psi e^{i\Delta}; ~~~~ \abs{\psi'}^2 = \abs{\psi}^2
\end{equation*}

因此，如果波函数的相位移动任意一个常数，物理定律仍然不变，这是另一种对称性。物理定律必须具有这样的性质，即量子力学相位的移动不会产生什么差别。我们刚才说过，在量子力学中，对每个对称性都存在着一个守恒定律。与量子力学相位相关联的守恒定律看来是电荷守恒定律。总之，这是一件非常有趣的事情！

\section{镜面反射}

其次一个问题是在空间反射 下的对称性。本章余下的大部分篇幅都将用来讨论这件事。我们要问：物理定律在反射下是否对称？可以用以下方式说得更具体一些：假定我们制造了一件东西，比方说一只钟，它带有许多齿轮，还有指针和数字。这只钟滴答滴答地走着，工作着，而钟里面有着卷紧的发条。我们从镜子里来看这个钟。问题并不是镜子里的钟像 什么。但是让我们实际地制造出 另一个正好同前一个钟在镜子中的映像一样的钟；每当原来的钟中有一个右旋的螺旋，我们就在另一个钟的对应位置上安装一个左旋的螺旋：前一个钟的钟面上刻着“2”字，就在后一个钟的钟面上对应地刻上一个“\reflectbox{2}”字；如果前者的发条是这样卷紧的，那么后者就以正好相反的方式卷紧。当我们做完这一切之后，就有了两个物理上的钟，它们彼此之间的关系就是物体和它的镜像的关系，然而我们要强调一下，它们都是实际存在的、物质的钟。现在的问题是：如果发条上得一样紧，两个钟在同样的条件下开始走动，那么这两个钟是否会永远那样滴答走动，就像一对精确的物与像一样（这是一个物理问题，而不是哲学问题）？我们对物理定律的直觉将认为，它们会如此。

我们猜想，至少在这对钟的情况下，空间的反射体现了物理定律的一种对称性，如果我们把每件事情从“右”变到“左”，并保持其他条件不变，就不能说出有何差别。所以，我们暂时假定这是正确的。如果确是如此，那么就不可能用任何物理现象来区分“右”和“左”，就像例如不能用物理现象来定义物体的某一绝对速度一样。所以，用任何物理现象来绝对地定义我们所谓的与“左”相反的“右”是什么意思应当是不可能的，因为物理定律应该是对称的。

当然，自然界并不一定 是对称的。比如，利用所谓的地理学无疑能够定义“右边”。例如假定我们站在新奥尔良看芝加哥，佛罗里达就在我们的右边（只要我们站在地面上），所以，我们能用地理学来定义“右”和“左”。当然，任何系统中的实际状况并不一定具有我们所谈到的对称性。这里的问题在于，物理定律 是否对称。换句话说，如果有一个像地球那样的天体，但组成它的尘土是“左”旋的，而且有一个像我们这样的人站在像新奥尔良那样的位置观望像芝加哥那样的城市，但由于每件东西都正好反过来，所以佛罗里达就在左边了，那么这种情况是否违反物理定律 呢？显然，每件事都左右互换的话，看来似乎并非不可能，这并不违反物理定律。

另一个要点是我们的“右”的定义不应当与传统有关。区分左和右的一个简易方法是到机械零件商店去随意取一个螺丝。那么多数是拿到一个右旋螺纹——也就是说，并不一定是右旋螺纹，但得到一个右旋螺纹的机会要比得到一个左旋螺纹的机会多得多。这是个传统或习惯的问题，或者是偶然的结果，所以并不是一个基本定律的问题。我们可以意识到，人人都能着手制造左旋螺丝！

所以，我们必须设法找到从根本上来说包含着“右旋”的某些现象。我们讨论的另一种可能性是偏振光经过比如糖水这种溶液时其偏振面会发生旋转这一事实。正如我们在第33章中所看到的，比方说，在某种糖的溶液中，偏振面是向右旋转的。这就是定义“右旋”的一种方法，因为我们可以在水中溶解一些糖，于是偏振面就向右旋转。但是糖是从生物体——植物中取得的，如果我们用人造的糖试验时，我们会发现，偏振面并不旋转 ！但是假如先在这些不引起偏振面旋转的人造糖溶液中放进一些细菌（它们会吃去一些糖），然后再滤去细菌，我们就会发现仍有一些糖剩下来（差不多是先前的一半），这一回偏振面确实也转动了，但却以相反 的方向转动！这看来颇令人迷惑，不过却很容易加以解释。

我们举另一个例子：有一种物质叫做蛋白质，在所有生物中都含有它，它对生命来
\begin{wrapfigure}{l}{0.55\textwidth}
    \centering
    \includegraphics[width=0.53\textwidth]{Chapter52/（a）L-丙氨酸（左）；（b）D-丙氨酸（右）.png}
    \caption{\\（a）L-丙氨酸（左）；（b）D-丙氨酸（右）}
    \label{figure:（a）L-丙氨酸（左）；（b）D-丙氨酸（右）}
\end{wrapfigure}
说是十分重要的。蛋白质由氨基酸链组成。图\ref{figure:（a）L-丙氨酸（左）；（b）D-丙氨酸（右）}是一种从蛋白质中产生的氨基酸的模型，它称为丙氨酸。如果是从真实的生物体蛋白质中提取出来的丙氨酸，分子的排列就如图\ref{figure:（a）L-丙氨酸（左）；（b）D-丙氨酸（右）}（a）所示。另一方面，如果我们设法用二氧化碳、乙烷、氨等合成丙氨酸（我们能够 合成它，这并不是一个复杂的分子），就会发现所制成的丙氨酸含有等量的如图\ref{figure:（a）L-丙氨酸（左）；（b）D-丙氨酸（右）}（a），（b）所示的两种结构的分子！前一种从生物体中得到的分子称作L-丙氨酸 ，另一种化学成分相同（也就是有相同的原子，原子的关系也相同）的分子，与“左旋”L -丙氨酸分子相比，是“右旋”的，它称为D-丙氨酸 。有趣的是当我们在实验室内由简单的气体合成丙氨酸时，得到的是两类分子的等量混合物。然而，生物所利用的只是L-丙氨酸（这不完全正确，生物体中各处都有一些D-丙氨酸的特殊应用，但很少见。所有的蛋白质都只利用L-丙氨酸）。现在如果我们制备两类丙氨酸，并用这种混合物喂给某种喜欢“吃”或消耗丙氨酸的动物，那么它不能利用D-丙氨酸，只能利用L-丙氨酸。这就是我们在糖水中所碰到的事，细菌只吃有用的糖，而留下了“错误”的一类糖（左旋糖也是甜的，但与右旋糖的味道不同）。

所以，看来生命现象似乎能区分出“左”与“右”，或者化学能够这样做，因为两种分子在化学上是不同的。但实际上并非如此！就所能进行的物理测量来说，例如，对能量和化学反应速度的测量，等等，如果我们使其他每件事都互为镜像，那么两类分子起同样的作用。一类分子将使光向右偏转，另一类分子则将使光在通过同样数量的溶液时往左偏转得正好一样多。这样，就物理学而言，这两类氨基酸同样符合要求。就我们今天对事物的理解来说，依照薛定谔方程的基本原理，两类分子所显示的特性应当完全对应，这样，一类有右旋作用，另一类则有左旋作用。但是，在生命过程中，却只有一种方式起作用。

人们推测其理由如下。比如，我们不妨假设在某个时刻生命不知怎么处在这样一种状况下：某些生物中的所有的蛋白质都包含有左旋的氨基酸分子，所有的酶也是有倾向性的——生命体中的每种物质都是有倾向性的——因此也就是不对称的。这样，在消化酶将食物中的化合物变为另一种化合物时，有一类化合物对酶来说是“合适”的，另一类却是不合适的（就像灰姑娘和拖鞋的那个故事。只是现在我们所试的是“左脚”）。就我们所知道的来说，在原则上，比方说，我们可以造出这样一只青蛙，其中所有的分子都是反过来的，每件事都像是一只真实青蛙的左旋镜像，因而这就是一个左旋蛙。这个左旋蛙可以很正常地活动一些时间，但是它会发现找不到东西吃，因为如果它吞下一只苍蝇，它的消化酶不能起作用，组成苍蝇的是一类“错误”的氨基酸分子（要么我们给青蛙一只左旋蝇）。就我们现在所知，如果每件事都反过来的话，化学与生命过程将照样进行下去。

如果生命完全是一种物理与化学的现象，那么蛋白质都由同样的螺旋状的分子所组成这件事情就只有这样来理解了：最初某一时刻，由于偶然的因素突然出现了某些生命分子，其中有的得到繁衍。在某个地方，一次有一个有机分子带有一定的倾向性，而从这一特殊事件出发，“右旋”就刚好在我们的特定环境下发展起来。一个个别的偶然历史事件是有倾向性的，但从这以后倾向性本身就传播开来。一旦达到今天的这种状态，当然它就将一直持续存在下去——所有的酶只消化“右旋”的东西，制造右旋的东西：当二氧化碳与水蒸气等等进入植物的叶中，制造糖的酶就把它们变成右旋的，因为酶本身就是右旋的。如果以后会产生什么新的病毒或活体的话，那么除非它们能“吃”已经存在的一类生命物质，不然就不能生存下去。因此它们也必须是同一类的东西。

这里不存在什么右旋分子数量上的守恒。右旋分子一旦突然出现，其数量就能保持继续增加。所以，人们推测，生命现象这种情况本身并不表明物理定律的缺乏对称性，相反，只是表明了宇宙的本性以及在上述含义下的地球上一切生命本源的共同性。

\section{极矢量与轴矢量}

现在我们进一步讨论下去。我们可以看到，物理学中许多地方都有着“右手”和“左手”规则。事实上，在学习矢量分析时，我们学到了必须用右手规则来正确地得出角动量、力矩、磁场等等的方向。例如，在磁场中运动的电荷所受的力就是$\boldsymbol{F} = q\boldsymbol{v}\times\boldsymbol{B}$。在给定的情况下，我们知道了$\boldsymbol{F}$、$\boldsymbol{v}$、$\boldsymbol{B}$，这个方程是否足以定义右旋性？实际上，如果我们回过去考虑一下矢量的来源，就知道“右手规则”只是一种习惯而已，它只是一种巧妙的方法。其实像角动量、角速度之类的量本来根本不是矢量！它们都是以某种方式与一定的平面相联系，只是因为空间有三维，所以可把有关量与垂直于那个平面的方向联系起来，在两种可能的取向中，我们选取了“右旋”的方向。

所以，如果物理定律是对称的，我们将会发现，如果某个魔鬼偷偷溜进所有的实验室，在每本有右手规则的书里用“左”这个词来替换“右”，因而我们一概使用“左手规则”的话，那么在物理定律上不会造成任何差别。

我们来作一点说明。矢量可以分为两类，有一类是“真正”的矢量，比如空间中的位移$\Delta\boldsymbol{r}$。如果在我们的仪器中，这里有一个零件，那里有另外一个零件，那么在一个镜像仪器里，有前一个零件的镜像物，也有后一个零件的镜像物。如果我们从“这个零件”到“那个零件”画出矢量，那么一个矢量就是另一个矢量的镜像（图\ref{figure:空间的位移矢量与其镜像}）。矢量箭头变换了指向，就好像整个空间翻了个身一样，这一种矢量我们称为极矢量。

\begin{figure}[htbp]
    \centering
    \begin{minipage}[t]{0.42\textwidth}
        \centering
        \includegraphics[width=6cm]{Chapter52/空间的位移矢量与其镜像}
        \caption{空间的位移矢量与其镜像}
        \label{figure:空间的位移矢量与其镜像}
    \end{minipage}
    \quad\quad
    \begin{minipage}[t]{0.42\textwidth}
        \centering
        \includegraphics[width=6.2cm]{Chapter52/转轮与其镜像。注意角速度“矢量”方向并没有反转}    
        \caption{转轮与其镜像。注意角速度“矢量”方向并没有反转}
        \label{figure:转轮与其镜像。注意角速度“矢量”方向并没有反转}
    \end{minipage}
\end{figure}

但是，另一类与转动有关的矢量具有不同的性质。例如，在三维空间中有某个物体在作转动，如图\ref{figure:转轮与其镜像。注意角速度“矢量”方向并没有反转}所示。如果在镜子中看它，将作如图右边所示的转动，也就是说作为原来那个转动的镜像而转动着。现在我们约定用同样的规则表示镜像的转动，它也是一个“矢量”，在反射后，并没有 像极矢量那样改变，但是相对于极矢量以及空间的几何关系而言，则正好反过来；这种矢量称为轴矢量。

现在，如果反射对称定律在物理上是正确的，我们必须这样来设计方程，即当我们改变每个轴矢量的符号和每个矢积的符号时（它相当于反射），不应出现任何差别。比如，当我们写出一个公式表明角动量为$\boldsymbol{L}=\boldsymbol{r}\times\boldsymbol{p}$ 时，这个方程是完全正确的，因为如果我们换成左手坐标系时，$\boldsymbol{L}$的符号改变了，而$\boldsymbol{p}$和$\boldsymbol{r}$没有改变；但矢积的符号变化了，因为我们要从右手规则变到左手规则。再举个例子，我们知道作用于在磁场中运动的电荷
\begin{wrapfigure}{r}{0.4\textwidth}
    \centering
    \includegraphics[width=0.38\textwidth]{Chapter52/磁铁与其镜像}
    \caption{磁铁与其镜像}
    \label{figure:磁铁与其镜像}
\end{wrapfigure}
上的力为$\boldsymbol{F} = q\boldsymbol{v}\times\boldsymbol{B}$，但当我们从右旋变到左旋系统时，由于$\boldsymbol{F}$和$\boldsymbol{v}$都是极矢量，所以由矢积所要求的变号应当被$\boldsymbol{B}$的变号所抵消，这就意味着$\boldsymbol{B}$必须是轴矢量。换句话说，如果进行这样一种反射，$\boldsymbol{B}$必须成为$-\boldsymbol{B}$。所以，在把坐标系从右手改为左手后，我们也必须使磁铁的两极互换。

我们用例子说明上述情况。假定我们有如图\ref{figure:磁铁与其镜像}所示的两块磁铁。一块磁铁上的线圈按某种方式缠绕，电流按一个确定的方向流过线圈。另一块磁铁就像前一块的镜像一样，线圈按相反的方式缠绕，在线圈内发生的每件事都正好反过来，电流方向如图所示。现在，按产生磁场的定律（这一点我们还没有正规地学习过，但多半在高中已知道一些），这里磁场的方向应如图中所示。若一块磁铁的一个磁极是南极，则在另一块磁铁中，电流按相反方向流动，这样磁场就反了过来，相应地出现一个北极。这样我们看到，从右旋改为左旋时，我们的确要把磁铁的南北极互换！

我们对磁极的改变不必介意，这些也都只是习惯而已。让我们谈谈现象 吧。假如现在有一个电子穿过一个磁场，进入纸面。于是，如果我们用公式$\boldsymbol{v}\times\boldsymbol{B}$来求电子所受的力$\boldsymbol{F}$（记住电荷是负的），就可发现，按照物理定律，电子将在确定的方向上发生偏转。这样物理现象就是，当一个线圈通以一定指向的电流时，电子的运动轨迹也按一种确定的方式弯曲——这就是物理内容——这里毋需考虑如何给每件事情贴上标记。

现在我们用一面镜子来做同样的实验：使电子通过与原来对应的方向，于是力的方向反了过来，如果我们按同样的规则计算它的话，结果是很好的，因为对应的运动 是一种镜像运动！

\section{哪一只是右手}

实际的情况是，在研究任何现象时，总有两个或偶数个右手规则，而最后结果是：现象看起来总是对称的。因此，简言之，如果我们不能区别南极与北极，那么也无法分别左与右。然而看来我们好像可以 说出磁铁的北极。例如，罗盘指针的北极就是指向北方。但实际上这也是一种与地理学有关的局部特征，这正像有关芝加哥所在方向的谈论一样，是不算数的。如果我们见过罗盘指针，就会注意到，指北针是浅蓝色的。但这是人们涂到小磁针上去的颜色。这些都是局部性的、约定的判定标准。

但是，如果磁铁真的具有这种性质：当我们十分靠近地观察它时，就能看到在北极而不是在南极上长有细丝，如果这是一般的规律，或者如果存在着任何 其他独特的区分磁铁南北极的方法，我们就能说出实际情况是两种可能情况中的哪一种，而这就是反射对称定律的终结 。

为了更清楚地说明整个问题，不妨设想我们同某个火星人，或某个极其遥远的理性生物通过无线电话进行交谈。我们不得发送给他任何实际的样品，以供观察，比如，假设我们能发送光信号给他，就可以送去右旋圆偏振光，并对他说：“这就是右旋光——你只要注意它的旋转方向就知道了。”但是，我们不能给 他送去任何东西，只能和他交谈。他离这里太远，或者在某一个奇怪的地方，以致不能看到任何我们能见到的东西。比如，我们不能说：“看一看大熊星座；请注意这些星是如何排列的。我们所指的‘右’是……”我们只能通过无线电话交谈。

现在要告诉他有关我们的所有事情。当然，首先我们要从数的定义开始，于是说：“滴答、滴答，二 ，滴答、滴答、滴答，三 ，……”这样他渐渐地能够理解几个词，等等。不一会儿我们就可能跟这个伙伴变得十分熟悉，于是他说：“你这个家伙究竟是什么样子？”我就开始自我描写，并且告诉他说：“噢，我们有6 ft高。”他便说：“等一等，6 ft是什么意思？”有没有可能告诉他6 ft是多长吗？当然行！我们可以说：“你知道氢原子的直径吧！我们有17000000000个氢原子那么高！”这之所以可能，是因为物理定律在尺度变化时不是不变的，因而我们可以 定义一个绝对长度。这样我们就解释了自己身材的尺寸，并且把我们的一般形状也向火星人作了描述——有四肢，在四肢上有五个手指或脚趾，等等，他也就顺着我们来进行想象。我们描述自己的外形时，料想不会遇到任何特殊的困难。在我们讲述的过程中，火星人甚至还做了一个有关我们的外形的模型。接着他就说：“哎呀！你真是个非常漂亮的家伙，但是在你的身体内有些什么呢？”于是我们就开始描写身体内的各种器官，随后，我们谈到心脏，在仔细地描写了心脏的形状之后，我们就说：“现在请把心脏的位置安排在左边。”他就问道：“慢着，什么是左边？”于是我们的问题就是向他描写心脏在哪一边，而他既不能看到我们所见到的任何东西，我们也不能向他发送任何我们所谓的“右”的样品——没有一个标准的右旋的物体。我们能这样做吗？

\section{宇称不守恒}

我们知道，万有引力定律、电磁定律、核力都符合反射对称原则，所以，这些定律以及任何由它们推得的东西都不能应用。但是，与自然界中发现的许多基本粒子相关联，
\begin{wrapfigure}{l}{0.45\textwidth}
    \centering
    \includegraphics[width=0.43\textwidth]{Chapter52/粒子衰变的简图}
    \caption{$\tau^+$粒子衰变和$\theta^+$粒子衰变的简图}
    \label{figure:粒子衰变的简图}
\end{wrapfigure}
存在着一种称为$\beta$衰变 或弱衰变 的现象。其中弱衰变的一个例子与大约在1954年发现的粒子有关，它使人们感到很难理解。有一种带电粒子蜕变为三个$\pi$介子，如图\ref{figure:粒子衰变的简图}所示。这个粒子一度称为$\tau$介子。在图\ref{figure:粒子衰变的简图}中我们还看到另一个曾称为$\theta$介子的粒子蜕变成两个介子。根据电荷守恒，其中一个必须是中性的。这样，一方面我们有一个称为$\tau$的粒子蜕变成三个$\pi$介子，还有一种$\theta$粒子则蜕变成两个$\pi$介子。不久人们发现$\tau$与$\theta$在质量上几乎相等。事实上，在实验误差之内，它们是相等的。其次，人们发现，它们分别蜕变为三个$\pi$介子与两个$\pi$介子所需的时间也几乎相等；并且具有相同的寿命。再有，无论何时生成这两种粒子时，它们总以同样的比例出现，比如说，14\%是$\tau$介子，86\%是$\theta$介子。

任何头脑清楚的人都立即认识到：它们必定是相同的粒子，我们只是产生了一个有两种不同蜕变方式的东西，而不是两种不同的粒子。所以，这个可以按两种方式蜕变的东西具有同样的寿命和同样的产品比例（因为这就是粒子进行两类不同蜕变的可能性的比例）。

但是，可以证明（我们在这里完全无法说明如何 证明），根据量子力学中的反射对称原理，不可能 由同一种粒子得到这两种结果——同一个粒子不可能 以两种这样的方式蜕变。与反射对称原理相对应的守恒定律没有经典的类比，这一类量子力学的守恒关系称作宇称守恒 。这样，由于宇称守恒，或更确切地说，由于弱衰变的量子力学方程对反射的对称性，同一种粒子不可能按两种方式变化，所以这必定是某种质量、寿命等等方面的巧合。但是，人们越深入研究，这种巧合也愈加惊人，于是人们逐渐产生了疑问：也许，深奥的自然界反射对称定律可能并不正确。

由于出现了这种明显的失败，李政道和杨振宁建议做一些有关衰变的其他实验，试图检查一下定律在其他情况下是否正确。第一个这样的实验是由哥伦比亚大学的吴健雄女士做的，她的实验如下：人们知道，蜕变时发射一个电子的钴的某一同位素处于极低温度和极强的磁场中时是磁化的，如果温度低到使热振动不致过分扰动原子磁体的话，这些原子磁体就在磁场中排列起来。所以，钴原子在这个强磁场中就全都排列起来。随后它们发射一个电子而蜕变，人们发现，当原子排列在$\boldsymbol{B}$矢量朝上的磁场中时，大多数的电子是在朝下的方向上发射出去的。

如果一个人不是真正地“熟悉”世界，那么这种议论似乎没有丝毫意义。但是如果他懂得世界上的问题和有趣事情的话，就会看出这是一个最为戏剧性的发现：当我们把钴原子放到极强的磁场中去时，向下的蜕变电子比向上的蜕变电子要多。所以如果我们在“镜子”中进行对应的实验，钴原子将沿着相反的方向排列，此时它们就将往上 而不是往下 发射较多的电子，情况是非对称的 。磁体长出细丝了 ！磁铁的南极成为这样一种磁极：在β衰变中电子趋向于离开它！于是，这就在物理上区别了南北极。

在这以后，人们还做了其他许多实验：$\pi$介子蜕变为$\mu$和$\nu$介子；$\mu$介子蜕变为一个电子和两个中微子；近来，$\Lambda$蜕变为质子和$\pi$，$\Sigma$的蜕变以及许多其他的蜕变实验。事实上，在几乎所有可以预期的情况中，全都发现不 遵从反射对称原则！从根本上说，在物理学的这一个阶梯上，反射对称定律是不正确的。

简言之，我们能够告诉火星人该把心脏放到哪一个部位了。我们可以这样说：“听着，自己制造一块磁铁，把线圈绕上去，让电流通过，随后取一些钴，并使温度降低，然后再这样来安排实验，使电子从脚部向头部运动，那么电流通过线圈时，流进的方向就是我们称之为右的方向，而流出的方向就是左的方向。”所以，现在只要做一个这样的实验，就能够确定右与左了。

人们还曾预言过许多其他特征。例如，已知钴核的自旋，即角动量在蜕变前是$5\hbar$，在蜕变后是$4\hbar$。电子带有自旋角动量，还牵涉到中微子。从这里很容易看出电子必须具有与其运动方向一致的自旋角动量，中微子也同样如此。所以看上去好像电子往左自旋，这也得到实验验证。事实上，博姆和瓦帕斯特拉就是在我们这里验证了电子大多数是向左旋转的（还有另一些实验给出了相反的结果，但它们是错误的）。

另一个问题当然就是要找出宇称守恒失败的规律。有没有什么法则能告诉我们这种不守恒的情况在多大的范围内成立？有。这个法则是，只有在非常慢的称为弱衰变的反应中，守恒才遭到破坏，而且在这种情况发生时，有关的法则表明，带有自旋的粒子，例如电子、中微子等等，在出现时倾向于向左自旋。这是一条倾向一面的法则，它把速度极矢量与角动量轴矢量联系起来，并且指出角动量与速度方向相反的可能性比一致的可能性要来得大一些。

这条法则就是如此，但今天我们并不真正理解它的原因。为什么这条法则是正确的，它的基本原因是什么，它与其他事情有何联系？这件非对称的事实使我们感到如此的震惊，以致此刻还没有能从惊讶中充分地恢复过来去理解对于所有其他规则来说这将意味着什么。然而，这个课题是有趣和新颖的，也是仍未获得解决的，所以看来我们讨论一些与此有关的问题是可取的。

\section{反物质}

当一种对称性丢失之后，我们要做的第一件事就是赶快查一下已知的或假定成立的对称性的表，看看是否还会失去什么别的对称性。在我们的表中没有提到一种对称操作，这就是物质与反物质的关系，它也必须受到质疑。狄拉克曾预言除电子外必定还有另一种称为正电子的粒子（由安德森在本学院发现），它必然与电子有关。这两种粒子的所有性质都服从一定的对应法则：能量相等；质量相等；电荷相反。但是，比所有其他都重要的是，当它们碰在一起时，就彼此湮没而把所有的质量转化为能量，例如，以$\gamma$射线的形式释放掉。正电子称为电子的反粒子 ，而这些就是粒子与它的反粒子的特征。从狄拉克的论证可以清楚地看出，世界上所有其余的粒子也应有对应的反粒子。比如，对质子应有反质子，用符号$\bar{p}$来表示。$\bar{p}$应具有负电荷，它的质量和质子相同，等等。然而最重要的特点是，质子和反质子碰在一起彼此就会湮没。我们强调这件事的原因在于，人们对我们所说的有中子也有反中子这一点感到不好理解，他们说：“中子就是中性的，那么又怎么可能 有相反的电荷呢？”“反”粒子的规则并不只是说它具有相反的电荷，它有一系列特征，所有这些特征都是相反的。反中子和中子的区别就在于，如果我们把两个中子放在一起，它们仍然是两个中子，但是如果把一个中子和一个反中子放在一起，它们彼此会湮没，并且释放出巨大的能量，发射出各种$\pi$介子、$\gamma$射线，等等。

现在，如果我们有了反中子、反质子和反电子，在原则上就可以造出反原子。虽然到现在还没有造出反原子，但在原则上这是可能的。例如，一个氢原子在中心有一个质子，在外面有一个电子绕着转动。现在设想在某个地方我们能产生一个反质子，而在外面带着一个正电子，正电子会不会绕着转动？会。首先，反质子带负电，而反电子带正电，这样它们就以相应的方式互相吸引——正负粒子的质量是一样的，每件事都一样。这是物理学的对称原理之一，方程式似乎表明，如果我们用某种物质制造出一个钟，然后又用反物质制造出一个同样的钟，它将同样走动（当然，如果把两个钟放在一起，它们就都会湮没，但那是另一回事了）。

这样就出现了一个问题。我们可以用物质制造出两个钟，一个是“左旋”的，另一个是“右旋”的。例如，我们可以不用简单的方式制造钟，而使用钴和磁铁以及电子探测器（它能检测出$\beta$衰变电子的存在，并对之计数）来制造钟。每计数一次，秒针就走动一下。那么另一个镜像钟由于接收到较少的电子，将不会走得一样快。所以，显然我们可以制造两个钟：一个左旋，另一个右旋，它们走得不一样。那就让我们用物质制造一个可称为标准的或“右旋”的钟，另外还有这种物质制造一个左旋的钟。我们刚才已看出，一般说来，这两者不 会走得一样快，而在那个著名的物理现象发现之前，人们曾认为它们会走得一样快。我们还假定正物质与反物质是等价的。那就是说，如果我们用反物质制造一个同样形状的右旋钟，那它将会走得和右旋正物质钟一样快，而如果制造同样的左旋钟，它也会走得一样快。换句话说，原先人们相信所有这四个钟 都是相同的，现在我们当然已经知道右旋物质和左旋物质并不一样。因此，可以假设，右旋反物质和左旋反物质也并不一样。

一个明显的问题是，究竟哪两种钟是相同的，如果发生这种情况的话？换句话说，右旋物质钟与右旋反物质钟走得一样快吗？或者说，右旋物质钟与左旋反物质钟走得一样快吗？利用正电子衰变来代替电子衰变的$\beta$衰变实验指出了这里的相互关系是：右旋物质的行为与左旋反物质的行为一样。

于是，现在终于可以说右与左的对称性仍然保持着！如果我们用反物质代替正物质制造一个左旋钟，它将走得一样快。这样事情就变为，代替我们的对称性表中的两条独立规则的，是把这两者结合在一起变成一条新规则，即右旋的物质与左旋的反物质是对称的。

这样，如果火星人是由反物质造成的，而我们若指点他如何作一个类似于我们的“右旋”模型，当然，结果就刚好相反。在我们之间进行了许多交谈后，我们彼此互相教会对方制造一艘宇宙飞船，然后乘飞船在空间半途相遇，那么会发生什么事呢？我们会把彼此的传统和习惯等等告诉对方，并且大家会很快地跑过去情不自禁地伸出手来。要是他真的伸出的是左手，请千万小心！

\section{对称破缺}

其次一个问题是，我们怎么来理解接近于 对称的定律？令人惊异的是，一方面，在物理学的很大范围内，包括核力的强作用现象，电磁现象，以及最弱的引力现象等重要领域中，关于这些现象的所有定律似乎都是对称的。另一方面，那个小小的例外部分则跑出来说：“不，定律并不都是对称的！”自然界几乎对称，但又不完全对称，这究竟是怎么一回事？我们怎样来理解这一点？首先，我们是否还有什么别的例子？答案是，事实上我们确实有一些别的例子。比如，在质子与质子之间，中子与中子之间，中子与质子之间，相互作用的核力部分都完全相同，这里有着一种核力的对称性，一种新的对称性，所以，我们可以交换中子与质子——但很明显，这并不是普遍成立的对称性，因为两个相隔一定距离的质子间的电斥力对中子来说并不存在。所以，一般而言，并不总是 能用中子来代换质子，这种代换只是一个良好的近似。为什么是良好的 近似？因为核力远远比电力强，所以这也是一种“几乎”对称的情况。这样，我们在别的事情上确实也看到了例子。

在我们的心目中有一种倾向，认为对称是无比完美的。事实上，这与希腊人的一个古老观念相类似：圆是完美的，如果去相信行星的轨道不是圆形，而只是接近于圆形的话，这就太可怕了。是一个圆和近似于一个圆这两件事之间的差别不是一个很小的差别，对于我们的认识来说，这是一种根本性的改变。在圆上存在着对称性与完美性的迹象，一旦稍有偏离，就一切都完了，它就不再有对称性。于是，问题在于为什么行星的轨道只是接近于 圆——这是一个更加困难的问题。一般地说，行星的实际运动轨道应当是椭圆形的。但是在漫长的岁月里，由于潮汐力的作用等等因素，这些轨道变得几乎对称了。现在的问题是，我们这里是否也存在着类似的事情。若从圆的观点出发来看，如果轨道都是精确的圆，这种情况显然很简单，也自然用不到去解释。但是既然轨道只是接近于圆，就需要作许多解释，结果表明这是一个很大的动力学问题，于是，我们就得考虑潮汐力等等的影响来解释为什么轨道是近于对称的。

这样，我们的问题就是要解释对称性究竟从何而来。为什么自然界是如此近于对称？没有人能道出所以然。我们可能作出的唯一解释大致如此：日本的日光市有一座门，这座门有时被日本人称为全日本最美的城门，它是在深受中国艺术影响的时代建造的。这个城门非常精致，有许多山墙和美丽的雕刻图案，还有许多柱子以及刻有龙头及贵族雕像的圆柱，等等。但是当你挨近看时，在一根柱子上除了见到复杂精细的雕刻图案外，还可见到有个小小的图样刻得正好颠倒过来。要是没有这件事，情况就完全对称了。如果你要问为什么会那样，据说有这样一个传说：它是故意刻得颠倒的，为的是使上帝不致妒忌人的完美。人们故意在这里留下一个小小的错误，那样上帝就不会因为妒忌而对人类感到愤怒了。

我们愿意把这种看法反过来，并且相信自然界之所以接近于对称，其真正的解释是：上帝只将物理定律造得接近于对称，这样我们就不会妒忌上帝的完美了！